\documentclass[11pt]{article}
\usepackage[utf8]{inputenc}
\usepackage{ctex}
\usepackage{amsmath, amssymb, amsthm}
\usepackage{geometry}
\geometry{a4paper, margin=1in}

\input{../preamable/preamable.tex}

\declaretheorem[style=thmgreenbox, name=定义]{cndefinition}
\declaretheorem[style=thmredbox, name=定理]{cntheorem}
\declaretheorem[style=thmbluebox, name=性质]{cnproperty}
\declaretheorem[style=thmredbox, name=引理]{cnlemma}
\declaretheorem[style=thmredbox, numbered=no, name=推论]{cncorollary}
\declaretheorem[style=thmbluebox, numbered=no, name=例]{cnexample}

\title{近世代数笔记(Lec 06)：Sylow定理与有限阿贝尔群的结构}
\author{Raysance}
\date{}

\begin{document}
\maketitle

\section{Sylow 定理}

在拉格朗日定理中，我们知道有限群子群的阶必定整除群的阶，但是反过来，如果 $d \mid |G|$，是否一定存在阶为 $d$ 的子群？Sylow 定理对这种逆问题给出了在特定素数幂情况下的肯定回答。

\begin{cntheorem}[Sylow 第一定理]
设 $G$ 是有限群，且 $|G|=p^n m$，其中 $p$ 是素数，$n \ge 1$，且 $\gcd(p,m)=1$。则 $G$ 一定存在阶为 $p^n$ 的子群。

通常，我们将这样的阶为群阶中包含 $p$ 的最高次幂的子群称为 \textbf{Sylow $p$-子群}。
\end{cntheorem}

\begin{proof}
我们可以通过群作用来证明此定理。

令集合 $\mathcal{S}$ 为 $G$ 中所有大小为 $p^n$ 的子集构成的集合，即：
\[ \mathcal{S} = \{ S \subset G \mid |S| = p^n \} \]
其势为组合数 $|\mathcal{S}| = \binom{p^n m}{p^n}$。可以证明，该组合数不能被 $p$ 整除（即 $|\mathcal{S}| \not\equiv 0 \pmod p$）。

定义群 $G$ 在集合 $\mathcal{S}$ 上的左乘作用：对于 $g \in G, S \in \mathcal{S}$，定义 $g \cdot S = \{ gs \mid s \in S \}$。
由于 $G$ 作用在 $\mathcal{S}$ 上，$\mathcal{S}$ 可以划分为若干个轨道。因为 $|\mathcal{S}|$ 不是 $p$ 的倍数，所以必定存在至少一个轨道 $C$，它的长度 $|C|$ 也不是 $p$ 的倍数。

取该轨道中的一个元素 $S \in C$，令 $H = \text{Stab}(S)$ 是 $S$ 在 $G$ 中的稳定子群，即 $H = \{ g \in G \mid gS = S \}$。
根据轨道-稳定子定理：
\[ |C| = \frac{|G|}{|H|} = \frac{p^n m}{|H|} \]
因为 $|C|$ 不被 $p$ 整除，故 $p^n$ 必须整除 $|H|$，这意味着 $|H| \ge p^n$。

另一方面，对于任意 $s \in S$，由于 $H S = S$，必然有 $H s \subseteq S$，从而 $|H| = |Hs| \le |S| = p^n$。
综合以上两点结果，我们得出 \textbf{$|H| = p^n$}。
因此，有限群 $G$ 必然存在阶为 $p^n$ 的子群 $H$。
\end{proof}

\begin{cntheorem}[Sylow 第二定理]
设 $G$ 是有限群，且 $|G|=p^n m$（$p \nmid m$）。则 $G$ 中任意两个 Sylow $p$-子群都是共轭的。而且，$G$ 的任何 $p$-子群都包含在某个 Sylow $p$-子群中。
\end{cntheorem}

\begin{proof}[直观理解与部分证明]
设 $P$ 和 $Q$ 都是 $G$ 的 Sylow $p$-子群。考虑 $Q$ 在 $G/P$ (即 $P$ 的全体左陪集组成的集合 $\mathcal{S} = \{ g_1 P, \dots, g_m P \}$，注意这不一定是商群，因为 $P$ 不一定是正规子群) 上的左乘群作用：
对于 $q \in Q$，$q \cdot (g_iP) = (qg_i)P$。

这里 $|\mathcal{S}| = [G:P] = m$，而 $p \nmid m$。因为 $Q$ 的阶是 $p$ 的幂次，根据 $p$-群作用的性质，必定存在不动点。
即存在一个陪集 $gP$，使得 \mbox{$\forall q \in Q$，$qgP = gP$}，即 $g^{-1}qgP = P$，这就意味着 $g^{-1}Qg \subseteq P$。
但 $|Q| = |P| = p^n$，因而 $g^{-1}Qg = P$，也就是 $Q$ 和 $P$ 是共轭的。
\end{proof}

\begin{cntheorem}[Sylow 第三定理]
设 $G$ 是有限群，且 $|G|=p^n m$（$p \nmid m$）。令 $r_p$ 为 $G$ 中所有 Sylow $p$-子群的个数，则有：
\[ r_p \mid m \quad \text{并且} \quad r_p \equiv 1 \pmod p \]
\end{cntheorem}

直观理解：Sylow 第三定理给出了在寻找或者限制一定阶数群的 Sylow $p$-子群数量的极强约束条件。利用前两条性质：$r_p$ 其实等于正规化子群的指数指数或由共轭类的性质推导出。

\section{直和与循环群}

\begin{cntheorem}[循环群的直和性质]
设 $\mathbb{Z}_m$ 和 $\mathbb{Z}_n$ 分别是阶为 $m$ 和 $n$ 的循环群，则直和 $\mathbb{Z}_m \oplus \mathbb{Z}_n$ 也是循环群（且同构于 $\mathbb{Z}_{mn}$），当且仅当 $\gcd(m,n)=1$。
\end{cntheorem}

性质理解：如果 $\gcd(m,n) > 1$，则所有元素的阶的最大值是 $\text{lcm}(m,n) < mn$，此时不可能生成包含 $mn$ 个元素的整个群。这也是下面将有限阿贝尔群继续分解的基础。

\section{有限阿贝尔群的结构定理}

有限阿贝尔群具有极其良好的结构，可以使用一组唯一的数（或其直和分解）来完全刻画任意有限阿贝尔群的同构类。

\begin{cntheorem}[有限阿贝尔群的基本定理]
任何有限阿贝尔群都同构于若干个循环群（其阶均为素数的幂次）的直和。 

设 $|G| = n = p_1^{k_1} p_2^{k_2} \cdots p_m^{k_m}$ 是群阶的素因数分解。则存在唯一的一组正整数 $a_{ij} \ge 1$，其中针对每个 $i \in \{1,2,\dots,m\}$ 有：
\[ k_i = a_{i1} + a_{i2} + \dots + a_{it_i} \]
使得 $G$ 同构于：
\[ G \cong \left( \mathbb{Z}_{p_1^{a_{11}}} \oplus \dots \oplus \mathbb{Z}_{p_1^{a_{1t_1}}} \right) \oplus \dots \oplus \left( \mathbb{Z}_{p_m^{a_{m1}}} \oplus \dots \oplus \mathbb{Z}_{p_m^{a_{mt_m}}} \right) \]
这组数序列 $\big( p_1^{a_{11}}, p_1^{a_{12}}, \dots, p_m^{a_{mt_m}} \big)$ 被称为 $G$ 的 \textbf{初等因子组 (Elementary Divisors)}。
\end{cntheorem}

利用该结构定理，我们可以在给定群的阶时，分类出所有可能的有限阿贝尔群的同构类。

\begin{cnexample}
求所有 36 阶的阿贝尔群的同构类。
\end{cnexample}
\textbf{解}：我们将群阶 36 进行素因子分解：$36 = 2^2 \times 3^2$。

素因子 2 的指数为 2，它对应的所有拆分（Partition）有：
\begin{enumerate}
    \item $2 = 2$ ，此时对应 $\mathbb{Z}_{2^2} = \mathbb{Z}_4$
    \item $2 = 1+1$ ，此时对应 $\mathbb{Z}_{2^1} \oplus \mathbb{Z}_{2^1} = \mathbb{Z}_2 \oplus \mathbb{Z}_2$
\end{enumerate}

素因子 3 的指数也为 2，它对应的所有拆分有：
\begin{enumerate}
    \item $2 = 2$ ，此时对应 $\mathbb{Z}_{3^2} = \mathbb{Z}_9$
    \item $2 = 1+1$ ，此时对应 $\mathbb{Z}_{3^1} \oplus \mathbb{Z}_{3^1} = \mathbb{Z}_3 \oplus \mathbb{Z}_3$
\end{enumerate}

将上面独立的情况组合，总共有 $2 \times 2 = 4$ 种不同的群结构组合（在同构意义下）：
\begin{enumerate}
    \item $\mathbb{Z}_4 \oplus \mathbb{Z}_9 \cong \mathbb{Z}_{36}$ ——(对应的初等因子为 4 和 9)
    \item $\mathbb{Z}_2 \oplus \mathbb{Z}_2 \oplus \mathbb{Z}_9 \cong \mathbb{Z}_2 \oplus \mathbb{Z}_{18}$ ——(对应的初等因子为 2, 2, 9)
    \item $\mathbb{Z}_4 \oplus \mathbb{Z}_3 \oplus \mathbb{Z}_3 \cong \mathbb{Z}_3 \oplus \mathbb{Z}_{12}$ ——(对应的初等因子为 4, 3, 3)
    \item $\mathbb{Z}_2 \oplus \mathbb{Z}_2 \oplus \mathbb{Z}_3 \oplus \mathbb{Z}_3 \cong \mathbb{Z}_6 \oplus \mathbb{Z}_6$ ——(对应的初等因子为 2, 2, 3, 3)
\end{enumerate}
上述这四种结构刻画了所有可能的 36 阶有限阿贝尔群。

\end{document}
